\section{Open problems}\label{ch:fmopen}

The last two problems concern the variable $e^{2X}$ of Section~\ref{ch:fmexp}, where no hypothesis on the zeros enters.

\begin{problem}\label{fm:prob:shoulder}
Mixed shoulders. Theorem~\ref{fm:thm:shoulder} decides the symmetric laws on $[-1,1]$ whose density is concave on $[0,1]$ and those whose density is convex there and not linear on its support: by \eqref{fm:eq:sine} the question is whether the sine transform of the slope $g=-p'$ has only real zeros, and the theorem answers it when $g$ is monotone. Decide the slopes that first rise and then fall. One class is settled. If $g$ is symmetric about $\frac12$, then $p(0)=\int_0^1g=1$, because $1=2\int_0^1p=2\int_0^1xg(x)\,dx=\int_0^1g$ by the symmetry of $g$, so $u\mapsto g(\frac12+u)$ is a probability density on $[-\frac12,\frac12]$; and $X\overset d=\frac12U+B$, with $B$ independent of $U$ and distributed with that density, since the density of $\frac12U+B$ has slope $g$ on $(0,1)$. As $\sin(z/2)/(z/2)$ has only real zeros, $X$ is then ferromagnetic exactly when $B$ is. For the beta slopes $g(y)\propto y^a(1-y)^b$, $a,b\ge0$, this decides the diagonal $a=b$, where $B$ is a symmetric beta law and $X$ is ferromagnetic (Section~\ref{fm:sec:examples}), and Theorem~\ref{fm:thm:shoulder} decides the edges: $X$ is ferromagnetic for $b=0$, where $p$ is concave, and not for $a=0<b$, where $p$ is convex and not linear. The chapter does not decide the rest of the quadrant. On the half line the slope of the P\'olya law rises and then falls (Section~\ref{fm:sec:polyamix}), and RH is the statement that its sine transform has only real zeros.
\end{problem}

\begin{problem}\label{fm:prob:gs}
The Griffiths Simon class. Newman's theorem places the P\'olya law in the class of Ellis, Monroe and Newman (Section~\ref{fm:sec:polyanec}). Does it belong to the Griffiths Simon class of Section~\ref{fm:sec:def}? The laws of that class are ferromagnetic, so a positive answer would imply RH, and under Ingham's hypothesis the spins of any approximating system would have to be coupled (Section~\ref{fm:sec:nosqrt}).
\end{problem}

\begin{problem}\label{fm:prob:split}
Splitting the zeros. Assume RH, so that $\varphi_X(t)=\prod_\gamma(1-t^2/\gamma^2)$ over the positive ordinates, with no Gaussian factor by Theorem~\ref{fm:thm:nosqrt}. Can the ordinates be split into two nonempty sets $A$ and $A'$ such that $\prod_{\gamma\in A}(1-t^2/\gamma^2)$ and $\prod_{\gamma\in A'}(1-t^2/\gamma^2)$ are both characteristic functions? This is the same as writing $X=D+R$ with $D$ and $R$ independent, ferromagnetic and not $0$, and by Theorem~\ref{fm:thm:nosqrt} both sets are then infinite. Under Ingham's hypothesis no three ordinates are in arithmetic progression, so no piece is a scaled fair sign or a scaled uniform variable, whose zeros form arithmetic progressions.
\end{problem}

\begin{problem}\label{fm:prob:zc}
Positivity of the zero completed laws. The laws $f_K$ of Section~\ref{fm:sec:zerocompleted} are proved to be ferromagnetic probability densities for $K\le16$ by the Maxwell construction, and for every $K\le200$ and for $K=300$, $400$, $600$ and $800$ by a computer assisted proof in ball arithmetic. Compute the threshold $v^*(K)$ near $K=5000$, where the linear and the power law fits of the ratio $v^*(K)/2\beta_K$ part; the cancellation there calls for a working precision of tens of thousands of bits. Is $v^*(K)<2\beta_K$ for every $K$? If so, every $f_K$ is a ferromagnetic probability density by Theorem~\ref{fm:thm:zc}(ii), and with no hypothesis their limit in law, with characteristic function $e^{-\beta_\infty\theta^2}\prod_k(1-\theta^2/\gamma_k^2)$ over the zeros on the critical line, is a ferromagnetic law by Theorem~\ref{fm:thm:closed}; by Theorem~\ref{fm:thm:zc}(iii) it is the P\'olya law exactly when RH holds.
\end{problem}

\begin{problem}\label{fm:prob:omega}
The cosine constant. By Section~\ref{fm:sec:tilts}, $\Omega\le\frac12$ with no hypothesis, $\Omega\le\frac38$ if the half plane $\re s>\frac78$ announced in October 2026 is free of zeros (Section~\ref{sec:ch25:open}), and RH holds if and only if $\Omega=0$. Is the set of $a\ge0$ for which $\cosh(at)\,p$ is ferromagnetic an interval? Improve these bounds by computation, as computations with the zeros have improved de Bruijn's bound $\Lambda\le\frac12$ on the de Bruijn Newman constant, to $\Lambda\le0.2$ \cite{PT21}. A first case is $a=\sqrt{0.1}$, below both bounds: the law $\cosh(\sqrt{0.1}\,t)\,p$ is as close to $p$ as the flow at $\lambda=0.2$, which is ferromagnetic, and deciding it is a computation with the zeros of $\frac12(\xi(s+a)+\xi(s-a))$, of the kind that gave the bound $0.2$ for the flow.
\end{problem}

\begin{problem}\label{fm:prob:sbp}
The size bias principle for real shapes. Theorem~\ref{fm:thm:sbp} is proved for integer shapes. Does it hold for all real shapes $t_n>0$? For $K=1$ it is the gamma identity, with $G=0$. Integer shapes enter the proof twice, through the divided difference formula for the density of $Y$ and through Lemma~\ref{fm:lem:sign}, and together they show that on each gap $(\tau_k,\tau_{k+1})$ the phase of the Laplace transform of $Y_q$ is nonincreasing and falls by exactly $q$. Is this still true for real shapes? With it, the rest of the proof goes through, provided the estimates of the transform near the points $-\tau_n$ and at infinity can be checked as well.
\end{problem}

\begin{problem}\label{fm:prob:powers}
The powers of $e^X$. For $2<|\beta|$ below the HCM threshold, which lies between $5.5791395$ and $\pi/0.563096=5.57914219\ldots$ (Proposition~\ref{fm:prop:hcm}), is $e^{\beta X}$ infinitely divisible, self decomposable or a GGC? Below the threshold it has no HCM density, which leaves the GGC property open, and Theorem~\ref{fm:thm:sd} gives self decomposability only at $|\beta|=2$. Is the HCM threshold exactly $\pi/v_m=5.5791396$? That requires $h\ge0$ on the open strip $0<v<v_m$, where $h(0,v)\to0$ as $v\uparrow v_m$, so a certificate must replace the balls with positive lower bounds near $iv_m$ by a local argument.
\end{problem}
